% Canonical geometric integration of the 2026-10-09 H2 compactification study.
The punctured models of Section~\ref{sec:p0-holed} can be organized by their
ends without identifying a singular point, a boundary circle, and a computational
input.  We first specify which structure is compactified, then construct a
model for every finite positive puncture count, and finally state the collar
data needed to cut and reconstruct the original AES.

The results concern existence and reconstruction under explicit hypotheses.
They neither classify all AES metrics nor identify the complete hyperbolic
metric of an auxiliary uniformization with the arithmetic metric.  The port
binding and execution layer is developed in Paper~IV
\cite{Yuan2026AEGCondensation}; its inputs do not alter the geometry defined here.

\subsection{End data and three meanings of completion}
\label{subsec:p0-end-data}

\begin{definition}[Finite-type carrier and conformal point compactification]
\label{def:p0-finite-type-carrier}
A finite-type carrier of genus $\gamma$ with $k$ punctures is a connected
oriented surface $M$ diffeomorphic to $\overline M\setminus S$, where
$\overline M$ is a closed oriented surface of genus $\gamma$ and $|S|=k$.
It is a punctured AES carrier when $(g,a;\mu,\lambda)$ also satisfies
Definition~\ref{def:p0-punctured-aes}, including essentiality at every point of
$S$.  Its oriented conformal structure is the one determined by $g$.
A \emph{conformal point compactification} adds one point at each end and
extends that conformal structure; it does not assert extension of $g$ or $a$.
\end{definition}

\begin{proposition}[The local condition for conformal filling]
\label{prop:p0-conformal-end-filling}
A finite-type conformal surface admits such a point compactification if and
only if every end has a neighbourhood conformally equivalent to a punctured
disc $0<|z|<\epsilon$.  A topological puncture alone does not impose this
condition.
\end{proposition}

\begin{proof}
If a conformal compactification exists, a disc coordinate at each added point
restricts to the required end coordinate.  Conversely, attach $z=0$ in each
of finitely many disjoint end charts.  Their overlap maps are already
holomorphic away from the added points, and no other chart contains those
points; hence these charts extend the conformal atlas.  The finite-type
compact core together with the finitely many closed smaller end discs is
compact.  A finite-modulus annular end, for example, is not conformally a
punctured-disc end although it is topologically annular.
\end{proof}

For a complete hyperbolic cusp, the quotient of $\{\operatorname{Im}z>Y\}$
by $z\mapsto z+1$ has coordinate $q=e^{2\pi iz}$.  Direct substitution gives
\[
 h=\frac{|dq|^2}{|q|^2(\log|q|)^2}.
\]
Thus $q=0$ fills the conformal puncture, while the cusp remains at infinite
$h$-distance.  Topological point filling, conformal point filling, and metric
completion are different operations.  Removing an open disc instead leaves a
boundary circle; an oriented real blow-up can retain directions at a puncture,
but does not automatically extend the arithmetic data.  Markings, end germs,
and any chosen angular or port data must be retained explicitly.

For a genus-zero complete finite-area hyperbolic surface with $k\ge3$ cusps,
the conformal compactification is a sphere with $k$ distinct marked points.
Normalizing three points to $0,1,\infty$ leaves $k-3$ complex parameters,
modulo any explicitly chosen label symmetries.  This is a statement about the
carrier's conformal moduli, not a count of AES assignments or metrics.

\subsection{An explicit family with every puncture count}

\begin{theorem}[Compactifiable AES with any positive finite puncture count]
\label{thm:p0-arbitrary-k-punctures}
Let $k\ge1$.  For $k>1$ choose distinct real numbers $r_1,\ldots,r_{k-1}$;
for $k=1$ use the empty product.  Choose a polynomial $P_k$ with
\[
 P_k'(x)=\prod_{j=1}^{k-1}(x-r_j),\qquad
 a(x,y)=P_k(x)+y^2,
\]
and set $S=\{r_1,\ldots,r_{k-1},\infty\}\subset\PP^1(\CC)$,
$M=\PP^1(\CC)\setminus S$.  For $\mu\ne0$ and real $\lambda$, put
\begin{equation}
 g=\frac{(P_k'(x))^2+4y^2}{\mu^2+\lambda^2a^2}(dx^2+dy^2).
 \label{eq:p0-arbitrary-k-metric}
\end{equation}
Then $(\PP^1(\CC),S,g,a;\mu,\lambda)$ is a punctured AES with exactly
$k$ essential punctures and conformal compactification $\PP^1(\CC)$.
\end{theorem}

\begin{proof}
The Euclidean differential of $a$ is $(P_k'(x),2y)$, whose zeros are precisely
$(r_j,0)$.  The numerator in \eqref{eq:p0-arbitrary-k-metric} is consequently
positive on $M$ and the denominator is at least $\mu^2>0$.  Inverting its
positive conformal factor gives $|da|_g^2=\mu^2+\lambda^2a^2$.
At a finite point of $S$ the polynomial assignment has vanishing differential;
any smooth non-degenerate extension of the given AES would therefore contradict
that identity.  At infinity, along $x=0$ and $|y|\to\infty$, the assignment
$P_k(0)+y^2$ diverges, so it admits no continuous real-valued extension.
These prove essentiality directly.  The conformal class on $M$ is the
Euclidean class, which extends in the standard coordinate $1/z$ at infinity
and at all finite punctures.  For $k=1$, $P_1'=1$ and the only puncture is
infinity.  No critical-set theorem at infinity is used.
\end{proof}

This proves existence for every $k$, including the low-complexity cases
$k=1,2$; it does not supply every marked conformal structure or a prescribed
zero set.  The notation $\mathfrak E_k$ remains descriptive, not a unique
model determined by $k$.  In the four-puncture calibration we take
\begin{equation}
 a=\frac{x^4}{4}-\frac{x^2}{2}+y^2,
 \qquad S=\{-1,0,1,\infty\},\qquad \mu=\lambda=1.
 \label{eq:p0-k4-calibration}
\end{equation}
The finite puncture values are $-1/4,0,-1/4$.  They are fixed by this assignment;
arbitrary later computational inputs cannot be interpreted as arbitrary end
limits without constructing a different parameterized AES family.

\subsection{Three-ended blocks and exact geometric reconstruction}

For $2\gamma-2+k>0$, a finite-type carrier has a pants decomposition into
three-ended blocks.  An external end can remain a puncture, while an internal
end of a cut block is a boundary circle with a collar.  If $N$ is the number
of blocks and $E$ the number of internal seams, Euler characteristic and end
counting give
\begin{equation}
 N=2\gamma-2+k,\qquad 3N=k+2E,\qquad E=3\gamma-3+k.
 \label{eq:p0-pants-counts}
\end{equation}
Indeed each block has Euler characteristic $-1$, and gluing circles changes
no Euler characteristic.  Existence of the decomposition and its hyperbolic
coordinate interpretation are standard; see \cite[\S2]{Wolpert2003Completion}.
In genus zero the adjacency graph is a tree with $k-2$ vertices and $k-3$
internal edges.  If instead $k$ counts extra operational ends in addition to
three permanent ends, substitute $k+3$ in these formulas.  The exceptional
sphere, plane, and cylinder cases are not assigned negative numbers of blocks.

\begin{definition}[AES collar contract]
\label{def:p0-collar-contract}
A collar contract specifies the paired seam sides, their parametrizations and
markings, and smooth two-sided collar charts for the proposed quotient.
On overlapping collar charts, transition diffeomorphisms $\phi_e$ preserve the
surface orientation and satisfy
\begin{equation}
 \phi_e^*g_2=g_1,\qquad a_2\circ\phi_e=a_1,
 \label{eq:p0-collar-matching}
\end{equation}
with identical signed constants $\mu,\lambda$.  Inverse and cocycle identities
are required wherever transitions overlap.  Boundary orientations of paired
cut circles are opposite.  The contract includes the complete seam inventory,
external end labels and germs, and the requirement that the quotient atlas is
a Hausdorff smooth oriented surface.
\end{definition}

The conditions concern full collar germs, not just equal values on the seam.
Equivalently, one may specify compatible smooth tensors in one signed collar
coordinate and restrict them to its two sides.  Arbitrary boundary
identifications do not automatically satisfy the definition.

\begin{proposition}[Cutting and descent of the original AES]
\label{prop:p0-aes-collar-descent}
Cut a finite-type regular AES along a pants decomposition and retain its full
collar contracts and external end data.  Reassembly reconstructs the original
AES up to the recorded marked isomorphism.  More generally, AES blocks with
contracts as in Definition~\ref{def:p0-collar-contract} descend to a regular
AES on the smooth quotient.  Their canonical arithmetic frames agree across
the seams.  A finite motion trajectory crossing the seams transversely away
from punctures is carried to the same trajectory in the quotient.
\end{proposition}

\begin{proof}
The matching identities glue the metric tensors to a smooth positive-definite
metric and the assignments to a smooth function on the quotient atlas.  The
eikonal identity is local and invariant under change of chart, so it holds on
the quotient.  Orientation-preserving isometries preserving $a$ preserve its
gradient and the rotation $J$; with the same signed $\mu,\lambda$, the formula
in Proposition~\ref{prop:canonical-arithmetic-frame} therefore matches both
frame fields.  Uniqueness for smooth local ordinary differential equations
matches their flows on overlaps, hence each finite admissible transverse
crossing.  In the case obtained by cutting an existing AES, the quotient charts
are its original charts and the seam identifications restore exactly the same
points, tensors and function.  Retained end germs restore the same essential
punctures; no extension across them is asserted.
\end{proof}

For \eqref{eq:p0-k4-calibration}, the circles
$|z+1/2|=3/4$ and $|z-1/2|=3/4$ give two different single-seam decompositions.
The first separates $\{-1,0\}$ from $\{1,\infty\}$; the second separates
$\{0,1\}$ from $\{-1,\infty\}$.  Neither seam meets $S$.  Restricting the
same $a,g$ to either pair of blocks and retaining the original collar charts
reconstructs exactly the same four-puncture AES by the proposition.  This is a
reconstruction result for the recorded atlas, not independence from erased
seam data or a canonical choice of decomposition.

\subsection{Auxiliary hyperbolic geometry and three-cusp compatibility}

When the carrier admits a complete finite-area hyperbolic uniformization $h$,
its pants carry internal geodesic boundary lengths $\ell_e>0$ and twists
$\tau_e$.  With fixed marking and fixed external cusp types, these are
Fenchel--Nielsen coordinates on Teichm\"uller space.  They reconstruct $h$,
not the original AES tensors $(g,a)$.  On pinching every internal seam of a
chosen genus-zero decomposition, the limiting nodal curve consists of
$k-2$ thrice-punctured components, paired at $k-3$ nodes.  Local plumbing
$uv=q_e$ smooths a node when $q_e\ne0$; $q_e=0$ is its nodal limit
\cite[\S3]{Wolpert2003Completion}.  Local coordinates and any marking or
orbifold identifications must be specified.  This standard geometric
compatibility provides no AES deformation unless a compatible family
$(g_{\mathbf q},a_{\mathbf q})$ and its end data are also constructed.

\begin{warning}[The arithmetic metric is not the complete cusp metric]
\label{warn:p0-hyperbolic-auxiliary}
A connected, complete, boundaryless regular AES with a globally single-valued
assignment has underlying surface diffeomorphic to a plane or a cylinder.
Indeed the rectified assignment
$t(a)=\int_0^a(\mu^2+\lambda^2s^2)^{-1/2}\,ds$ has unit gradient; its complete
flow gives $M\cong a^{-1}(0)\times\mathbb R$, and a connected one-dimensional
level is a line or a circle.  This is the complete-splitting theorem of
Paper~I \cite{Yuan2026AEGFoundations}.  Consequently a complete three-cusp
hyperbolic metric cannot itself be the regular AES metric for such an
assignment.  Boundary, incompleteness, or a change of semantic category is an
active hypothesis, not a removable inconvenience.
\end{warning}

There is also an explicit special family covering the three-cusp sphere.
For $k\ge3$, put $d=k-2$ and
\[
 X_k=\PP^1(\CC)\setminus(\{0,\infty\}\cup\{z:z^d=1\}).
\]
The map $z\mapsto z^d$ is an unramified degree-$d$ cover from $X_k$ to
$\PP^1\setminus\{0,1,\infty\}$: any critical points lie in the removed set
$\{0,\infty\}$, and each retained target has exactly $d$ distinct preimages.
It extends to a branched map between the compact spheres.  This proves a
compatible conformal model for each count, not a covering for every prescribed
conformal structure.  For a four-puncture unramified cover the degree must be
$2$ by Euler characteristic; its compact extension has two simple branch
values among the three marked targets, and the four source punctures have
cross-ratio orbit $\{-1,2,1/2\}$.  Generic four-puncture moduli are excluded.
This covering construction by itself specifies no arithmetic metric.

\subsection{An assignment identity can retain motion residue}
\label{subsec:p0-motion-residue}

The scalar tearing of Section~\ref{sec:p0-acs-tearing} compares induced affine
operators.  A word inducing the identity on the assignment need not be the
identity on the full AES state.  Here is an exact local certificate in
\eqref{eq:p0-k4-calibration}, independent of a numerical trajectory plot.

\begin{proposition}[A nonzero level-set residue in the four-puncture model]
\label{prop:p0-k4-motion-residue}
Let $A_s$ and $M_t$ be the local flows of the canonical fields $X_u,X_v$ in
\eqref{eq:p0-k4-calibration}.  For sufficiently small $s,t$, the chronological
word $A_s,M_t,A_{-e^t s},M_{-t}$ defines
\[
 W_{s,t}=M_{-t}\circ A_{-e^t s}\circ M_t\circ A_s,
 \qquad a\circ W_{s,t}=a.
\]
At $p=(0,1)$ it has nonzero mixed displacement:
\[
 \left.\partial_s\partial_t W_{s,t}(p)\right|_{s=t=0}
 =([X_u,X_v]-X_u)(p)=(-3/4,0).
\]
Thus for sufficiently small nonzero $s=t$ the full state moves while its
assignment is unchanged.
\end{proposition}

\begin{proof}
The flow identities are $a\circ A_s=a+s$ and $a\circ M_t=e^t a$.
Their displayed composition gives the identity on $a$ exactly.  Write
$v=\nabla_{\mathrm{Eucl}}a$, $Jv=(-a_y,a_x)$, and $G=|v|^2$.  The canonical
fields are
\[
 X_u=(v-aJv)/G,\qquad X_v=(av+Jv)/G.
\]
At $p$, $a=1$, $v=(0,2)$ and $G=4$.  Differentiating the polynomial and these
rational fields gives
$[X_u,X_v](p)=(-1/4,1/2)$ and $X_u(p)=(1/2,1/2)$.
The usual two-parameter Taylor expansion of the four smooth local flows
has mixed term $st[X_u,X_v]$; the extra factor $e^t$ in the third flow
subtracts $stX_u$.  Since $W_{s,0}=W_{0,t}=\id$, its remainder along $s=t$
is $O(s^3)$, proving the nonzero displacement.  Finally
$da([X_u,X_v]-X_u)=\mu\lambda-\mu=1-1=0$,
so the infinitesimal residue is tangent to a level set, as required by the
exact assignment identity.
\end{proof}

An interface that retains only the endpoint number can therefore discard a
geometrically observable residue.  Whether that residue is needed for a
particular computation is a question about its declared future observations,
answered by Paper~IV's contextual residual criterion.  Neither pants counts,
circle packings, nor scalar assignment identities prove a faithful
realization of Adva's triadic computation.  The local circle surgery and
sheet constructions that motivated this programme, with their precise
non-equivalences, are recorded in Appendix~\ref{app:p0-circle-surgery}.
